\section{Introduction}
\label{sec:intro}
Automatic prompt and program optimisation has become routine: a tool such as DSPy~\citep{khattab2023dspy} or MIPRO~\citep{opsahlong2024mipro} searches instructions, demonstrations or multi-call pipelines on a set of benchmark items and the best configuration is reported with its score on benchmark items. When those items are the \emph{same} ones that drove the search, the number is a maximum over many noisy scores, not an unbiased estimate: it measures the winner's luck as well as the winner's quality. The consequence for the reader is unclear, because the size of the optimism, what governs it, and whether the comparisons built on such numbers (``our tuned system beats the baseline'', ``A beats B'') survive it are rarely examined for LLM systems.

We study this as a statistical question in a controlled setting. A search procedure with a budget evaluates $K$ configurations on $n$ items, returns $\hatc$, and the \emph{naive report} is the accuracy of $\hatc$ on those $n$ items; the target is its accuracy on fresh items. We ask (Q1) how large the gap is, what drives it, and whether it changes conclusions; and (Q2) whether more reliable reporting methods exist and how to test that they are.

\textbf{Contributions.} (1) A measurement of the optimism of the naive report in a synthetic world with exact truth (216 cells, 203{,}400 replicates) and on a real LLM and benchmark (Qwen2.5-1.5B-Instruct on 2000 MMLU-Pro questions, 300 prompt/ensemble configurations, 52 replay cells of 1000 replicates), with an extreme-value formula whose parameters are fitted in each setting (it checks the functional form; it is not a held-out prediction). (2) The decision-level consequence: naive reports declare a tuned system better than an untuned one when it is not, and rank two tuned systems by search effort and search-set size rather than by quality. These error rates rest on the synthetic world and, for real outcomes, on a paired A-versus-B design; the ``tuned beats baseline'' false-claim rate is not identifiable on MMLU-Pro, where the baseline is almost always truly worse. (3) A validation protocol with known truth and a real outcome matrix, applied to the reporting methods M0--M5 (naive, hold-out, cross-fitting, bootstrap optimism correction, simultaneous and conditional/hybrid selection-aware intervals, empirical-Bayes and Ladder-style disclosure), with an honest account of where each fails. (4) Scope: one small model on CPU, a prompt-format family rather than a search-heavy pipeline, emulated rather than live search (no DSPy/MIPRO run; the live tier and the 0.5B arm were not run), and a single fixed item pool whose between-pool variability is not estimated (Sec.~\ref{sec:discussion}); our real-LLM claims are limited accordingly.
